\section{Inquiry: restoration, direction, and fallibility}
\label{sec:inquiry}

Dissonance poses a question the current language cannot answer.  This
section characterizes what a successful answer must accomplish, shows that
the requirement directs inquiry without determining its output, and
establishes that any inquiry process meeting the paper's disciplines must be
fallible.

\subsection{Restorative refinements}

For $\mathcal P,\mathcal P'\in\Pi(S)$, write $\mathcal P\prec\mathcal P'$
when $\mathcal P'$ strictly refines $\mathcal P$.  After a refinement, the
retained record is recoded (\cref{ass:observation}) and the test is
re-applied with the recoded counts.

\begin{definition}[Restorative refinement]
\label{def:restorative}
$\mathcal P'$ is \emph{restorative} at $(\mathcal P_i,r_i)$ if
$\mathcal P_i\prec\mathcal P'$ and
$\widehat M_i(\mathcal P',r_i)\neq\emptysetset$ under the recoded record; it
is \emph{minimally restorative} if no restorative $\mathcal P''$ satisfies
$\mathcal P_i\prec\mathcal P''\prec\mathcal P'$.  Write
$\mathcal R(\mathcal P_i,r_i)$ for the set of minimally restorative
refinements.
\end{definition}

\begin{proposition}[Cellwise coherence]
\label{prop:coherence}
Call a cell of a refinement \emph{coherent} if some single mechanism passes
the test on the cell's recoded observations.  Then
$\widehat M_i(\mathcal P',r_i)\neq\emptysetset$ if and only if every cell of
$\mathcal P'$ is coherent.
\end{proposition}

\begin{proof}
Theories factor across cells---$M(\mathcal P')=\Theta^{\mathcal P'}$, and the
constraints in \cref{eq:test} apply cell by cell---so a passing assignment
exists exactly when a passing mechanism can be chosen for each cell
independently.
\end{proof}

Coherence is the general form of a conflict structure.  Token sets are
\emph{mutually incoherent} when no single mechanism accommodates their
union; every restorative refinement must separate mutually incoherent sets,
since a cell containing both would be incoherent.  In the deterministic
running example incoherence is a binary relation between labeled episodes and
restoration is a graph-coloring problem; in general, a mixture of three
mechanisms may defeat every single mechanism even though each pair of tokens
is individually accommodatable, so restoration is a hypergraph coloring.
Two structural facts hold in general: restoration is always possible (the
singleton partition recodes each token to its own state, where some mechanism
always fits), and---as the example shows vividly---minimal restorations are
typically multiple.

\subsection{The technology}

\begin{definition}[Representation-producing technology]
\label{def:technology}
A representation-producing technology for agent $i$ is a possibly randomized
map $\phi_i$ from $(\mathcal P_i,r_i)$ to
$\{\mathcal P_i\}\cup\mathcal R(\mathcal P_i,r_i)$: it fails, returning the
current partition, or succeeds, producing a minimally restorative
refinement.  It is \emph{truth-blind}: its output distribution depends only
on $(\mathcal P_i,r_i)$, never directly on $\tau$.  A complexity ordering
$\kappa$ on refinements---the analyst's model of which conceptual products
the agent's comparison capacities can generate---may restrict the range
further; $\kappa$ is a primitive of the technology, not an object of the
agent's optimization.
\end{definition}

The technology is the paper's stated representational primitive.  No formal
model derives a semantically new concept from nothing; the contribution is to
show how experience \emph{directs} a limited, fallible generative capacity
and why that capacity does not amount to a prior over future theories.  The
output, and even the set of possible outputs, is not a pre-insight
subjective menu: the agent cannot name a candidate refinement, assign it a
probability, or condition a policy on it before production succeeds
(\cref{sec:commitments}).

\subsection{Nonidentification and the price of robustness}

\begin{theorem}[Observational nonidentification]
\label{thm:nonidentification}
Let $\tau\neq\tau'$ be objective assignments inducing the same distribution
over agent $i$'s records under the committed profiles played.  Then any
truth-blind technology, and any truth-blind inquiry-and-formulation rule,
has the same output distribution in both worlds.  If no single theory is
$\eta$-adequate under both worlds at the visited pairs, no truth-blind rule
selects an adequate formulation with probability one in both.
\end{theorem}

\begin{proof}
The rule's input is the record, whose distribution is identical by
hypothesis, so its output distribution is identical.  If the adequate sets
are disjoint, any realized output is inadequate in at least one world with
the corresponding probability.
\end{proof}

\begin{proposition}[Robust sufficiency]
\label{prop:robust-sufficiency}
A partition sufficient for every assignment in a candidate set $\mathcal T$
must refine $\bigvee_{\tau'\in\mathcal T}\mathcal P^{\tau'}$, the join of the
candidates' causal-equivalence partitions.  A representation guaranteed to
accommodate all observationally equivalent worlds is accordingly finer than
any single world requires; a technology that economizes on distinctions must
accept fallibility.
\end{proposition}

\begin{proof}
Sufficiency for $\tau'$ is refinement of $\mathcal P^{\tau'}$
(\cref{def:awareness}); refinement of each candidate's partition is
refinement of their join.
\end{proof}

\begin{excont}
After the dissonant generation, the record contains two labeled episodes:
architecture $0$ objectively preferred at $s_{00}$, architecture $1$ at
$s_{11}$.  Of the fifteen partitions of four states, exactly four are
minimally restorative---each must separate $s_{00}$ from $s_{11}$: the row
partition $\{\{s_{00},s_{01}\},\{s_{10},s_{11}\}\}$, the column partition
$\{\{s_{00},s_{10}\},\{s_{01},s_{11}\}\}$, and the two partitions isolating
one labeled state.  The record directs inquiry (eleven partitions are
excluded) without determining it (four remain).  A natural complexity
ordering counts the coordinates needed to describe a refinement: rows and
columns are one-coordinate distinctions, the isolating partitions require
two.  A factor-seeking technology thus produces the row or the column
refinement---and here direction runs out.  Define the \emph{column world}
$\tau'(s_{rc})=\theta^{c}$: it agrees with the row world at both labeled
states, so the two worlds generate identical records
(\cref{thm:nonidentification} in exact form).  No truth-blind rule can
formulate correctly in both, a randomized rule succeeds in the worst case
with probability at most one-half, and by \cref{prop:robust-sufficiency} a
partition sufficient for both worlds must be the full partition into
singletons: robustness costs every distinction the language can draw, which
is exactly what a minimal technology economizes on.  The example also
locates the failure mode of each residual candidate: whichever factor
refinement arrives, its unique consistent theory extrapolates the wrong
prediction to the two unlabeled states in one of the two worlds.
\end{excont}
